\section{Supplementary: Alpha Ablation and Causal Sensitivity}

\label{app:alpha_ablation}

\subsection{Motivation}

A reviewer raised the concern that the paper's central claim --- that multiplicative
FiLM induces structural differentiation without functional coupling --- may be
contingent on the specific choice of modulation strength $\alpha = 0.1$. To address
this, we conduct a controlled sweep over $\alpha \in \{0.1, 0.5, 1.0\}$, measuring
the causal effect $\Delta = \mathcal{L}_{\text{trained}} - \mathcal{L}_{\text{removed}}$
at two time points: (1) immediately after loading the trained checkpoint (step~0), and
(2) after an additional 500 training steps at the target $\alpha$.

\subsection{Method}

Starting from the multiplicative FiLM checkpoint at step~850 (trained with $\alpha=0.1$,
$w_{\mathrm{std}} = 0.307$, scale projection weights frozen), we:

\begin{enumerate}
  \item Reload the checkpoint with $\alpha$ set to the target value.
  \item Run a single-pass causal evaluation on 10 held-out prompts (pure LM
    cross-entropy, no auxiliary loss terms).
  \item Train for 500 additional steps (same data, optimizer, and schedule as the
    original run) while keeping all other hyperparameters fixed.
  \item Repeat the causal evaluation.
\end{enumerate}

Each condition measures the same four outcomes: vanilla (no FiLM layers), trained
(normal forward), removed ($W_\ell$ active but input $S_t$ suppressed by setting
$\_s = \text{None}$), and frozen ($S_t$ fixed to $S_0$).

\subsection{Results}

\begin{table}[h]
\centering
\begin{tabular}{lcccc}
\toprule
$\boldsymbol{\alpha}$ & Step & $\Delta$ causal & $w_{\mathrm{std}}$ & Loss (trained) \\
\midrule
\multirow{2}{*}{0.1} & 0    & $-0.045$ & 0.307 & 2.700 \\
                     & 500  & $+0.095$ & 0.515 & 2.682 \\
\midrule
\multirow{2}{*}{0.5} & 0    & $+0.238$ & 0.307 & 2.938 \\
                     & 500  & $+0.526$ & 0.539 & 3.225 \\
\midrule
\multirow{2}{*}{1.0} & 0    & $+2.679$ & 0.307 & 5.381 \\
                     & 500  & $+2.862$ & 0.467 & 6.621 \\
\bottomrule
\end{tabular}
\caption{Alpha ablation results. $\Delta = \mathcal{L}_{\text{trained}} - \mathcal{L}_{\text{removed}}$;
positive values indicate that removing FiLM modulation \emph{improves} loss (modulation
is harmful). All loss values are pure LM cross-entropy on the same 10-prompt evaluation set.}
\end{table}

\subsection{Interpretation}

Three findings emerge:

\textbf{1. Causal sensitivity increases with $\alpha$.} The causal delta grows from
$-0.045$ at $\alpha=0.1$ to $+2.862$ at $\alpha=1.0$. This confirms that the
coupling strength between the modulation mechanism and the model's output is
controlled by $\alpha$ --- a trivial but necessary verification.

\textbf{2. Strong modulation is harmful, not helpful.} At $\alpha=0.5$ and $1.0$,
the causal delta is positive, meaning the trained model \emph{performs worse} with
FiLM enabled than without it. The modulation becomes destructive: the frozen LoRA
adapters and Qwen backbone were optimized under $\alpha=0.1$, and the abrupt change
in modulation amplitude disrupts their learned compensation patterns. Even 500
additional training steps fail to reverse this effect.

\textbf{3. The $\alpha=0.1$ regime is real but narrow.} The paper's central finding
--- structural differentiation without proportional functional influence --- holds
within the $\alpha=0.1$ operating point. However, the broader conclusion is that
effective modulation requires precise amplitude calibration: too small ($\alpha=0.1$)
produces measurable structure with marginal function; too large ($\alpha \geq 0.5$)
produces destructive interference. The narrow operating window ($\alpha \approx 0.1$)
represents a genuine architectural limitation of the current multiplicative FiLM
framework.<｜end▁of▁thinking｜>

\subsection{Implications for the Paper's Central Claim}

The reviewer's concern is addressed: the dissociation between $w_{\mathrm{std}}$ and
functional influence is confirmed at $\alpha=0.1$, the operating point used throughout
the paper. At higher $\alpha$, modulation becomes causally powerful but
destructive --- a different failure mode, not evidence of functional coupling.
Future work should explore learned adaptive $\alpha$ per layer, or gated routing
mechanisms that eliminate the need for fixed amplitude calibration.
